\sectionkicker{Source-backed technical notes}
\subsection*{総括――2024 H1に何が変わったのか: Technical Notes}
本欄は記事本文で圧縮した一次資料上の情報を、比較・再検証しやすい形で整理したものである。時系列、確認済みの主張、留意点、一次資料URLを掲載する。
\smallskip
\noindent{\small\textit{Technical Notesでは、一次資料から確認した公開・提供・研究上の事実を記載する。数値・能力評価や提供元・著者による比較表現は、独立再現済みの結果ではなく帰属claimとして扱う。}}\par
\medskip
\subsection*{Theme at a glance}
\addcontentsline{toc}{subsection}{Theme at a glance}
\begin{center}
\footnotesize
\begin{tabularx}{\linewidth}{@{}p{0.20\linewidth}p{0.10\linewidth}p{0.14\linewidth}X@{}}
\toprule
資料 & 位置づけ & 種別 & 時系列 \\
\midrule
GPT-4o & 補足資料 & モデル & 2024-05-13     \\
Meta Llama 3 & 補足資料 & オープンウェイト & 2024-04-18     \\
Apple Foundation Models & 補足資料 & モデル & 2024-06-10     \\
Jamba & 補足資料 & モデル & 2024-03-28     \\
Command R / Command R+ & 補足資料 & モデル更新 & 2024-03-11, 2024-04-04, 2024-03-14     \\
Sora & 補足資料 & モデル & 2024-02-15     \\
Anthropic interpretability research (Mapping the Mind / Golden Gate Claude) & 補足資料 & 公式情報 & 2024-05-23, 2024-05-21     \\
NVIDIA NIM & 補足資料 & Framework & 2024-03-18     \\
\bottomrule
\end{tabularx}
\end{center}
\normalsize

\begin{technicalnote}{GPT-4o}{補足資料}
\begingroup
% reader-facing Technical Notes break policy
\widowpenalty=10000
\clubpenalty=10000
\displaywidowpenalty=10000
\begin{tabularx}{\linewidth}{@{}>{\bfseries}p{0.22\linewidth}X@{}}
組織 & OpenAI \\
種別 & モデル \\
時系列 & 2024-05-13 \\
\end{tabularx}
\smallskip
{\bfseries 一次資料から整理したtechnical points}
\begin{itemize}[leftmargin=1.5em,itemsep=0.35em]
\item \textbf{一次情報で確認できる事実}: OpenAIの一次資料により、GPT-4oについて2024-05-13にモデル公開を確認した。 GPT-4oはtext・vision・audioを単一neural networkでend-to-end学習し、従来Voice Modeのspeech-to-text→LLM→text-to-speechという3-model pipelineでは失われていたtone・multiple speakers・background noise等を直接扱う設計として発表された。2024年5月13日のlaunch時点で公開されたのは主にtext/image inputとtext outputで、audio/video capabilityは段階提供予定だったため、model architecture上のmultimodalityと当日のproduct availabilityを分けて扱う。benchmark・latency・price比較はOpenAIによる評価として扱う。
\end{itemize}
\begin{minipage}{\linewidth}
% reader-facing Technical Notes generic boundary/source tail group
\begin{samepage}
{\bfseries 一次資料}
\begingroup\sloppy
\Urlmuskip=0mu plus 2mu\relax
\begin{itemize}[leftmargin=1.5em,itemsep=0.25em]
\item {\scriptsize\url{https://openai.com/index/hello-gpt-4o}}
\end{itemize}
\endgroup
\end{samepage}
\end{minipage}
\endgroup
\end{technicalnote}
\medskip

\begin{technicalnote}{Meta Llama 3}{補足資料}
\begingroup
% reader-facing Technical Notes break policy
\widowpenalty=10000
\clubpenalty=10000
\displaywidowpenalty=10000
\begin{tabularx}{\linewidth}{@{}>{\bfseries}p{0.22\linewidth}X@{}}
組織 & Meta \\
種別 & オープンウェイト \\
時系列 & 2024-04-18 \\
\end{tabularx}
\smallskip
{\bfseries 一次資料から整理したtechnical points}
\begin{itemize}[leftmargin=1.5em,itemsep=0.35em]
\item \textbf{一次情報で確認できる事実}: Metaの一次資料により、Meta Llama 3について2024-04-18にモデル公開を確認した。 Meta Llama 3の2024年4月公開分は8B/70Bのpretrained版とinstruction-fine-tuned版で、decoder-only Transformerを基盤に128K-token vocabularyのtokenizer、両サイズでgrouped query attention (GQA)、8,192-token sequence trainingを採用した。Metaはpretraining corpusを公開情報由来15T token超と説明しており、benchmark性能はMetaによる評価として扱う。
\end{itemize}
\begin{minipage}{\linewidth}
% reader-facing Technical Notes generic boundary/source tail group
\begin{samepage}
{\bfseries 一次資料}
\begingroup\sloppy
\Urlmuskip=0mu plus 2mu\relax
\begin{itemize}[leftmargin=1.5em,itemsep=0.25em]
\item {\scriptsize\url{https://ai.meta.com/blog/meta-llama-3/}}
\end{itemize}
\endgroup
\end{samepage}
\end{minipage}
\endgroup
\end{technicalnote}
\medskip

\begin{technicalnote}{Apple Foundation Models}{補足資料}
\begingroup
% reader-facing Technical Notes break policy
\widowpenalty=10000
\clubpenalty=10000
\displaywidowpenalty=10000
\begin{tabularx}{\linewidth}{@{}>{\bfseries}p{0.22\linewidth}X@{}}
組織 & Apple \\
種別 & モデル \\
時系列 & 2024-06-10 \\
\end{tabularx}
\smallskip
{\bfseries 一次資料から整理したtechnical points}
\begin{itemize}[leftmargin=1.5em,itemsep=0.35em]
\item \textbf{一次情報で確認できる事実}: Appleの一次資料により、Apple Foundation Modelsについて2024-06-10にResearch公開を確認した。 AppleがWWDC24で説明したfoundation-model構成は、約3B parametersのon-device language modelと、Apple silicon servers上のPrivate Cloud Computeで動くlarger server modelを組み合わせる。両モデルはgrouped-query attentionを使い、on-device側ではmixed 2-bit/4-bit LoRA-based palletizationを含む低bit最適化を適用するほか、featureごとのadapterをbase modelへ動的にload/swapしてtask specializationを行う。性能値はAppleによる評価として扱う。
\end{itemize}
\begin{minipage}{\linewidth}
% reader-facing Technical Notes generic boundary/source tail group
\begin{samepage}
{\bfseries 一次資料}
\begingroup\sloppy
\Urlmuskip=0mu plus 2mu\relax
\begin{itemize}[leftmargin=1.5em,itemsep=0.25em]
\item {\scriptsize\url{https://machinelearning.apple.com/research/introducing-apple-foundation-models}}
\end{itemize}
\endgroup
\end{samepage}
\end{minipage}
\endgroup
\end{technicalnote}
\medskip

\begin{technicalnote}{Jamba}{補足資料}
\begingroup
% reader-facing Technical Notes break policy
\widowpenalty=10000
\clubpenalty=10000
\displaywidowpenalty=10000
\begin{tabularx}{\linewidth}{@{}>{\bfseries}p{0.22\linewidth}X@{}}
組織 & AI21 Labs \\
種別 & モデル \\
時系列 & 2024-03-28 \\
\end{tabularx}
\smallskip
{\bfseries 一次資料から整理したtechnical points}
\begin{itemize}[leftmargin=1.5em,itemsep=0.35em]
\item \textbf{一次情報で確認できる事実}: AI21 Labsの一次資料により、Jambaについて2024-03-28にモデル発表を確認した。 JambaはTransformer、Mamba SSM、Mixture-of-Expertsを組み合わせたhybrid architectureで、AI21は52B total parametersのうちinference時に12Bをactiveにする構成と説明している。1/8層をTransformer layerとするblocks-and-layers設計、256K context、Apache 2.0のopen weightsが公開時点の主要仕様で、throughput/benchmark値はAI21による評価として扱う。
\end{itemize}
\begin{minipage}{\linewidth}
% reader-facing Technical Notes generic boundary/source tail group
\begin{samepage}
{\bfseries 一次資料}
\begingroup\sloppy
\Urlmuskip=0mu plus 2mu\relax
\begin{itemize}[leftmargin=1.5em,itemsep=0.25em]
\item {\scriptsize\url{https://www.ai21.com/blog/announcing-jamba/}}
\end{itemize}
\endgroup
\end{samepage}
\end{minipage}
\endgroup
\end{technicalnote}
\medskip

\begin{technicalnote}{Command R / Command R+}{補足資料}
\begingroup
% reader-facing Technical Notes break policy
\widowpenalty=10000
\clubpenalty=10000
\displaywidowpenalty=10000
\begin{tabularx}{\linewidth}{@{}>{\bfseries}p{0.22\linewidth}X@{}}
組織 & Cohere \\
種別 & モデル更新 \\
時系列 & 2024-03-11, 2024-04-04, 2024-03-14 \\
\end{tabularx}
\smallskip
{\bfseries 一次資料から整理したtechnical points}
\begin{itemize}[leftmargin=1.5em,itemsep=0.35em]
\item \textbf{一次情報で確認できる事実}: Cohereの一次資料により、Command R / Command R+について2024-03-11にモデル公開、2024-03-14にTool Use発表、2024-04-04にDeployment提供開始を確認した。 対象event近傍の一次資料から reranking / speculative decoding を確認できる。
\end{itemize}
\begin{minipage}{\linewidth}
% reader-facing Technical Notes generic boundary/source tail group
\begin{samepage}
{\bfseries 一次資料}
\begingroup\sloppy
\Urlmuskip=0mu plus 2mu\relax
\begin{itemize}[leftmargin=1.5em,itemsep=0.25em]
\item {\scriptsize\url{https://cohere.com/blog/command-r}}
\item {\scriptsize\url{https://cohere.com/blog/command-r-plus-microsoft-azure}}
\item {\scriptsize\url{https://cohere.com/blog/tool-use-with-command-r}}
\end{itemize}
\endgroup
\end{samepage}
\end{minipage}
\endgroup
\end{technicalnote}
\medskip

\begin{technicalnote}{Sora}{補足資料}
\begingroup
% reader-facing Technical Notes break policy
\widowpenalty=10000
\clubpenalty=10000
\displaywidowpenalty=10000
\begin{tabularx}{\linewidth}{@{}>{\bfseries}p{0.22\linewidth}X@{}}
組織 & OpenAI \\
種別 & モデル \\
時系列 & 2024-02-15 \\
\end{tabularx}
\smallskip
{\bfseries 一次資料から整理したtechnical points}
\begin{itemize}[leftmargin=1.5em,itemsep=0.35em]
\item \textbf{一次情報で確認できる事実}: OpenAIの一次資料により、Soraについて2024-02-15に研究Previewを確認した。 Soraの2024年2月technical reportは、可変duration・resolution・aspect ratioのvideo/imageを共同学習するtext-conditional diffusion modelとして記述し、時空間的に圧縮したvisual latentからspacetime patchesを切り出してTransformer tokenとして扱う。生成時はnoisy patchesからclean patchesを予測し、最大モデルでは最長1分のvideo生成を示したが、同資料はresearch previewとqualitative capability/limitation evaluationの範囲であり、後年の一般提供状態をH1へ遡及しない。
\end{itemize}
\begin{minipage}{\linewidth}
% reader-facing Technical Notes generic boundary/source tail group
\begin{samepage}
{\bfseries 一次資料}
\begingroup\sloppy
\Urlmuskip=0mu plus 2mu\relax
\begin{itemize}[leftmargin=1.5em,itemsep=0.25em]
\item {\scriptsize\url{https://openai.com/index/video-generation-models-as-world-simulators}}
\end{itemize}
\endgroup
\end{samepage}
\end{minipage}
\endgroup
\end{technicalnote}
\medskip

\begin{technicalnote}{Anthropic interpretability research (Mapping the Mind / Golden Gate Claude)}{補足資料}
\begingroup
% reader-facing Technical Notes break policy
\widowpenalty=10000
\clubpenalty=10000
\displaywidowpenalty=10000
\begin{tabularx}{\linewidth}{@{}>{\bfseries}p{0.22\linewidth}X@{}}
組織 & Anthropic \\
種別 & 公式情報 \\
時系列 & 2024-05-23, 2024-05-21 \\
\end{tabularx}
\smallskip
{\bfseries 一次資料から整理したtechnical points}
\begin{itemize}[leftmargin=1.5em,itemsep=0.35em]
\item \textbf{一次情報で確認できる事実}: Anthropicの一次資料により、Anthropic interpretability research (Mapping the Mind / Golden Gate Claude)について2024-05-21にResearch公開、2024-05-23にResearch公開を確認した。 Anthropicはdictionary learningをClaude 3 Sonnetのmiddle layerへscaleし、neuron activation patternsをhuman-interpretable featuresへ分解してmillions of featuresを抽出したと報告した。featuresはmultilingual/multimodalなconceptにも反応し、feature activationを人工的に増減するとmodel behaviorが変化することをGolden Gate Claude demoで示した。一方、抽出featureは全conceptの一部で、featureが属するcircuitsの理解は未解決と明記している。
\end{itemize}
\begin{minipage}{\linewidth}
% reader-facing Technical Notes generic boundary/source tail group
\begin{samepage}
{\bfseries 一次資料}
\begingroup\sloppy
\Urlmuskip=0mu plus 2mu\relax
\begin{itemize}[leftmargin=1.5em,itemsep=0.25em]
\item {\scriptsize\url{https://www.anthropic.com/news/golden-gate-claude}}
\item {\scriptsize\url{https://www.anthropic.com/research/mapping-mind-language-model}}
\end{itemize}
\endgroup
\end{samepage}
\end{minipage}
\endgroup
\end{technicalnote}
\medskip

\begin{technicalnote}{NVIDIA NIM}{補足資料}
\begingroup
% reader-facing Technical Notes break policy
\widowpenalty=10000
\clubpenalty=10000
\displaywidowpenalty=10000
\begin{tabularx}{\linewidth}{@{}>{\bfseries}p{0.22\linewidth}X@{}}
組織 & NVIDIA \\
種別 & Framework \\
時系列 & 2024-03-18 \\
\end{tabularx}
\smallskip
{\bfseries 一次資料から整理したtechnical points}
\begin{itemize}[leftmargin=1.5em,itemsep=0.35em]
\item \textbf{一次情報で確認できる事実}: NVIDIAの一次資料により、NVIDIA NIMについて2024-03-18に製品発表を確認した。 NVIDIA NIMはGTC 2024でinference microservicesとして提示され、NVIDIAのaccelerated-computing librariesとgenerative-AI modelsをパッケージ化し、industry-standard APIsで接続できる配備単位として説明された。NVIDIAはCUDA installed base上での利用、GPU世代に応じた再最適化、security vulnerability scanningを運用面の特徴として挙げている。
\end{itemize}
\begin{minipage}{\linewidth}
% reader-facing Technical Notes generic boundary/source tail group
\begin{samepage}
{\bfseries 一次資料}
\begingroup\sloppy
\Urlmuskip=0mu plus 2mu\relax
\begin{itemize}[leftmargin=1.5em,itemsep=0.25em]
\item {\scriptsize\url{https://blogs.nvidia.com/blog/2024-gtc-keynote/}}
\end{itemize}
\endgroup
\end{samepage}
\end{minipage}
\endgroup
\end{technicalnote}
\medskip
